\section{Research question, models, and evaluation criteria}
\label{sec:question}
\textbf{Objective.} Starting with random weights, train a model to predict the next music event and find the \emph{best performance within a fixed training-time budget per run}. This is a controlled study of model size, rather than an attempt to build a production music generator. The Smol Training Playbook recommends defining the purpose before choosing a model \citep{smolplaybook}; we use this practitioner guide for framing, not as peer-reviewed evidence for performance.

\textbf{Initial hypothesis.} A larger model will achieve lower validation loss because it can learn more complex musical patterns. We test this under the same two-hour GPU-time budget for each size. This is a hypothesis to check, rather than an expected result we must obtain. The initial comparison concerns model size; the follow-up question is developed from the baseline observations in Sections~\ref{sec:baselineanalysis}--\ref{sec:dropoutdesign}.

\begin{table}[H]
\centering
\caption{Dense decoder-only GPT variants. Counts include learned position embeddings and count tied token/output weights once.}
\begin{tabular}{lrrrr}
\toprule Model & Layers & Heads & Width & Total parameters\\\midrule
Small & 6 & 4 & 256 & 5,080,320\\
Medium & 8 & 8 & 512 & 25,891,328\\
Large & 12 & 12 & 768 & 86,029,056\\\bottomrule
\end{tabular}
\end{table}

Each run uses one Modal NVIDIA A10 GPU with 24\,GB memory and BF16 support \citep{nvidiaa10}; models run sequentially. All use vocabulary 376, context 1,024, and effective batch 32. Small/medium/large are project-specific labels, not pretrained GPT-2 variants.

\subsection*{Primary metric and checkpoint selection}
For target event $y_i$ and its earlier context $c_i$, the average negative log likelihood (cross-entropy loss) and perplexity are
\begin{equation}
\mathcal L(\theta)=-\frac{1}{N}\sum_{i=1}^{N}\log p_\theta(y_i\mid c_i),
\qquad \mathrm{PPL}=\exp\{\mathcal L(\theta)\}.
\label{eq:nll}
\end{equation}
Loss uses natural logarithms and is measured in nats/token. Lower loss and perplexity are better. Unlike accuracy, loss rewards assigning more probability to the correct event, even when it is not the top prediction. It gives no special credit for a musically similar wrong pitch and does not measure listening quality.

Each validation check averages 20 batches of four 1,024-token windows: 81,920 target positions, which may overlap. Resetting evaluation seeds gives all models the same windows. Dropout and gradient calculations are disabled.

A checkpoint is a saved model state. We select \code{best.pt}, the lowest recorded validation loss within the budget; \code{ckpt.pt} is the latest state. Selection uses no test loss. All runs have the full budget to improve, even if their best result occurs earlier. Final losses diagnose later overfitting. Section~\ref{sec:test} describes full-coverage evaluation.
